\section{Discussion}

\subsection{Overview of Main Findings}

Vehicle detection experiments in this study indicate that the final YOLOv11s model
provided the strongest performance for emissions-oriented vehicle counting from 
low-resolution traffic-camera imagery. Among the detector families evaluated, 
YOLO-based models consistently outperformed the Faster R-CNN and DETR alternatives.
The final selected YOLOv11s configuration achieved the strongest overall validation 
performance, reaching an AP50 of 0.857 and an mAP of 0.661, and it maintained strong
performance on the held-out test set with an AP50 of 0.772 and a mAP 0.607.

The most meaningful gains did not come from adopting a more complex architecture 
or from broad hyperparameter variation, but from targeted data-centric refinement.
In particular, revisiting failure cases through hard negative mining between 
passenger cars and passenger trucks produced a larger improvement than broader
oversampling strategies, local hyperparameter search, or continued fine-tuning 
from prior checkpoints. This suggests that, for this task, performance was constrained
less by a lack of architectural complexity than by ambiguity at specific class 
boundaries in the training data.

At the same time, the final detector generalized reasonably well beyond the validation 
set. Its held-out test-set performance remained broadly consistent with its validation 
results, indicating that the selected model generalized reasonably well to unseen 
images from the original dataset distribution. When evaluated on the separate 
cross-camera benchmark constructed from three unseen traffic cameras, performance
declined relative to the test set, but remained useful enough to support the broader
goal of deployment-oriented vehicle counting. Taken together, these results indicate
that traffic-camera imagery can support emissions-oriented vehicle detection and 
counting, despite the challenges introduced by low resolution, occlusion, and viewpoint 
shift.

% ==============================================================================

\subsection{Why the Final YOLO Pipeline Performed Best}



